\documentclass[12pt]{article}
\usepackage[utf8]{inputenc}
\usepackage[T1]{fontenc}
\usepackage{hyperref}
\title{Tarangan, East Language Resources for K-12 Curriculum Development}
\author{Compiled by Claude Sonnet 5 (claude-sonnet-5) for Andr\'as Kornai}
\date{}
\begin{document}
\maketitle
\section{Introduction}
East Tarangan (ISO 639-3 tre) is an Austronesian language of the Aru Islands, Maluku, Indonesia. It is one of the Aru languages, and the East and West varieties (tre and txn) differ enough that they may be no closer to each other than to Manombai. Speaker figures in the search results vary: about 14,000 for Tarangan as a whole (2011), and about 6,000 for East Tarangan in another source. Very little exists for this language, and I found no school materials.
\section{Primary Educational Resources}
I found no primers, textbooks or school programmes in East Tarangan.
\section{Linguistic and Reference Materials}
\begin{itemize}
\item Overview of the Tarangan language: \url{https://en.wikipedia.org/wiki/Tarangan_language}
\item Overview of the Aru languages: \url{https://en.wikipedia.org/wiki/Aru_languages}
\item East Tarangan numerals (MPI numerals database): \url{https://lingweb.eva.mpg.de/channumerals/Tarangan-East.htm}
\item Austronesian Comparative Dictionary language page (it cites an unpublished Hughes wordlist, which covers Dobel, not East Tarangan): \url{https://acd.clld.org/languages/581}
\item Glottolog reference entry for the Hughes wordlist: \url{https://glottolog.org/resource/reference/id/679273}
\end{itemize}
\section{Religious and Community Resources}
I found none.
\section{Digital and Audio Resources}
\begin{itemize}
\item Glosbe East Tarangan page. It is crowd-sourced, with unverified quality: \url{https://glosbe.com/tre/cbw}
\item PARADISEC field notes series for the Aru and Kei Islands, including historical Dutch vocabulary material on the Aru languages in general: \url{https://paradisec.org.au/fieldnotes/INIJAK.htm}
\item Endangered Languages Project entry: \url{https://endangeredlanguages.com/node/96515}
\end{itemize}
\section{Curriculum Development Context}
The language is not known to be taught in schools. No dictionary, grammar or teaching materials turned up in search. Curriculum work would have to start with new documentation, ideally with the community and with specialists in Aru languages.
\end{document}